\section{Introduction}
\label{sec:intro}
% 3D scene understanding
Understanding 3D scenes is essential for interacting with the environment in robotic navigation~\cite{cadena2016slam, mur_artal_2017_orbslam2}, autonomous driving~\cite{geiger2012kitti, sun2020waymo}, and augmented reality~\cite{klein2007ptam, izadi2011kinectfusion}. Traditional approaches, however, are constrained to a fixed set of object categories defined at training time~\cite{qi2017pointnet, choy2019minkowski, thomas2019kpconv}, limiting their applicability to open-world scenarios.
% open-vocabulary
Thanks to recent advances in vision-language models~\cite{radford2021clip, jia2021align}, open-vocabulary methods~\cite{vild2021, detic2022, openscene2023} enable querying and interacting with 3D scenes through natural language, and recognizing unseen object categories without requiring retraining.

While classical 3D understanding methods operate on point clouds or meshes derived from 3D sensors, recent neural scene representations, such as NeRFs~\cite{mildenhall2020nerf} and 3D Gaussian Splatting (3DGS)~\cite{kerbl20233d}, have emerged as a compelling alternative. They not only enable high-quality rendering from novel viewpoints but also facilitate semantic reasoning, as appearance and geometry are encoded jointly. Thus, open-vocabulary reasoning has recently been grounded in neural 3D scene representations~\cite{kerr2023lerf, qin2024langsplat}, enabling new semantic interactions in 3D environments.
Initially explored with NeRFs~\cite{kerr2023lerf,engelmann2024opennerf}, the efficiency and explicit nature of 3DGS simplified the integration of semantic features, contributing to its widespread adoption~\cite{qin2024langsplat,wu2024opengaussian,drsplat25,cheng2024occam}.


\begin{figure}[t]
    \centering
    % \fbox{\rule{0pt}{2in} \rule{0.9\linewidth}{0pt}}
   \includegraphics[width=\linewidth]{figs/teaser.pdf}
   \vspace{-5mm}
    \caption{Thanks to its feature aggregation that is visibility-aware and multi-view consistent, our proposed VALA is the most accurate and as quick as the fastest~\cite{cheng2024occam} to optimize. Comparison in 3D open-vocabulary segmentation on the LeRF-OVS dataset~\cite{qin2024langsplat}.}
    \vspace{-5mm}
    \label{fig:teaser}
\end{figure}

At the core of these approaches lies the challenge of embedding reliable semantic and language features into the 3D representation. Current methods rely on powerful off-the-shelf 2D foundation models, such as SAM~\cite{kirillov2023sam} and CLIP~\cite{radford2021clip}, which produce 2D feature maps that must be lifted to 3D and aggregated across views. Proper aggregation is critical for accurate 3D segmentation.

Despite numerous recent advances~\cite{drsplat25,jiang2025votesplat,li2024instancegaussian,sun2025cags}, current approaches suffer from an inherent limitation: they assign 2D features indiscriminately to \textit{all} Gaussians along a camera ray, disregarding scene geometry and occlusion relationships. Consequently, features originating from foreground objects (\eg, a vase) are incorrectly propagated to background structures (\eg, the supporting table or floor), leading to substantial degradation in open-vocabulary recognition accuracy.

Furthermore, when lifted into 3D, 2D features exhibit multi-view inconsistencies. The same object may produce divergent feature representations across different viewpoints, a phenomenon known as semantic drift~\cite{kerr2023lerf}. Current methods address this by promoting cross-view consistency through 3D-consistent clustering and contrastive objectives derived from SAM masks~\cite{wu2024opengaussian, liang2024supergseg, li2024instancegaussian, piekenbrinck2025opensplat3d}. Nevertheless, such strategies generally require extensive per-scene optimization, and their heavy reliance on noisy, view-dependent 2D cues often undermines cluster reliability.

In this paper, we address these fundamental feature aggregation problems with VALA (Visibility-Aware Language Aggregation), a lightweight yet effective framework that combines a two-stage gating mechanism with a robust multi-view feature aggregation strategy. Our gating mechanism leverages the statistical distribution of per-ray Gaussian contributions (termed visibility) to preferentially propagate features to Gaussians with high visibility, thereby ensuring accurate feature assignment. To further mitigate multi-view inconsistencies in 2D language features, we introduce a convex but non-smooth optimization on the unit hypersphere, which we reformulate into a streaming gradient-based procedure that achieves consistent embeddings without additional computational overhead. As shown in Figure~\ref{fig:teaser}, VALA strategies are highly effective.

Our contributions can be summarized as follows:
\begin{itemize}
    \item We identify fundamental issues in the feature aggregation of current works as a bottleneck in open-vocabulary 3D scene understanding.
    \item We introduce VALA, a visibility-aware feature propagation framework that employs a two-stage gating mechanism to assign features based on Gaussian visibility.
    % to tackle this problem with a visibility-aware feature propagation leveraging a two-stage gating mechanism.
    \item We propose a robust aggregation strategy for the 2D features using the streaming cosine median, thereby improving multi-view consistency.
    \item We obtain state-of-the-art performance in 2D \textit{and} 3D on open-vocabulary segmentation for 3DGS scenes on the reference datasets LeRF-OVS~\cite{qin2024langsplat} and ScanNet-v2~\cite{dai2017scannet}.
\end{itemize}
